\documentclass[10pt,a4paper]{amsart}
\usepackage{amsmath,amssymb,amsthm,graphicx}
% For commutative diagrams we use Xy-pic
\usepackage[all]{xy}

% We use 2cell for 2-commutative diagrams.
\xyoption{2cell}
% Theorem environments.
%
\theoremstyle{plain}
\newtheorem{theorem}[subsection]{Theorem}
\newtheorem{proposition}[subsection]{Proposition}
\newtheorem{lemma}[subsection]{Lemma}

\theoremstyle{definition}
\newtheorem{definition}[subsection]{Definition}
\newtheorem{example}[subsection]{Example}
\newtheorem{exercise}[subsection]{Exercise}
\newtheorem{situation}[subsection]{Situation}

\theoremstyle{remark}
\newtheorem{remark}[subsection]{Remark}
\newtheorem{remarks}[subsection]{Remarks}

\numberwithin{equation}{subsection}

% Macros
%
\def\lim{\mathop{\mathrm{lim}}\nolimits}
\def\colim{\mathop{\mathrm{colim}}\nolimits}
\def\Spec{\mathop{\mathrm{Spec}}}
\def\Hom{\mathop{\mathrm{Hom}}\nolimits}
\def\Ext{\mathop{\mathrm{Ext}}\nolimits}
\def\SheafHom{\mathop{\mathcal{H}\!\mathit{om}}\nolimits}
\def\SheafExt{\mathop{\mathcal{E}\!\mathit{xt}}\nolimits}
\def\Sch{\mathit{Sch}}
\def\Mor{\mathop{\mathrm{Mor}}\nolimits}
\def\Ob{\mathop{\mathrm{Ob}}\nolimits}
\def\Sh{\mathop{\mathit{Sh}}\nolimits}
\def\NL{\mathop{N\!L}\nolimits}
\def\CH{\mathop{\mathrm{CH}}\nolimits}
\def\proetale{{pro\text{-}\acute{e}tale}}
\def\etale{{\acute{e}tale}}
\def\QCoh{\mathit{QCoh}}
\def\Ker{\mathop{\mathrm{Ker}}}
\def\Im{\mathop{\mathrm{Im}}}
\def\Coker{\mathop{\mathrm{Coker}}}
\def\Coim{\mathop{\mathrm{Coim}}}

% Boxtimes
%
\DeclareMathSymbol{\boxtimes}{\mathbin}{AMSa}{"02}

%
% Macros for moduli stacks/spaces
%
\def\QCohstack{\mathcal{QC}\!\mathit{oh}}
\def\Cohstack{\mathcal{C}\!\mathit{oh}}
\def\Spacesstack{\mathcal{S}\!\mathit{paces}}
\def\Quotfunctor{\mathrm{Quot}}
\def\Hilbfunctor{\mathrm{Hilb}}
\def\Curvesstack{\mathcal{C}\!\mathit{urves}}
\def\Polarizedstack{\mathcal{P}\!\mathit{olarized}}
\def\Complexesstack{\mathcal{C}\!\mathit{omplexes}}
% \Pic is the operator that assigns to X its picard group, usage \Pic(X)
% \Picardstack_{X/B} denotes the Picard stack of X over B
% \Picardfunctor_{X/B} denotes the Picard functor of X over B
\def\Pic{\mathop{\mathrm{Pic}}\nolimits}
\def\Picardstack{\mathcal{P}\!\mathit{ic}}
\def\Picardfunctor{\mathrm{Pic}}
\def\Deformationcategory{\mathcal{D}\!\mathit{ef}}
\newenvironment{reference}{}{}
\newenvironment{slogan}{}{}
\newenvironment{history}{}{}
\usepackage{cleveref}
\crefname{originalequation}{original equation}{original equations}
\Crefname{originalequation}{Original equation}{Original equations}
\makeatletter
\newlabel{equation-form}{{0C65}{source}{}{sourcepointer-1}{}}
\newlabel{equation-form@cref}{{[originalequation][0][]0C65}{[1][0][]source}{}{}{}}
\makeatother
\crefname{originallemma}{original lemma}{original lemmas}
\Crefname{originallemma}{Original lemma}{Original lemmas}
\makeatletter
\newlabel{lemma-divisor-special-fiber}{{0C5Z}{source}{}{sourcepointer-2}{}}
\newlabel{lemma-divisor-special-fiber@cref}{{[originallemma][0][]0C5Z}{[1][0][]source}{}{}{}}
\makeatother
\crefname{originallemma}{original lemma}{original lemmas}
\Crefname{originallemma}{Original lemma}{Original lemmas}
\makeatletter
\newlabel{lemma-intersection-pairing}{{0C64}{source}{}{sourcepointer-3}{}}
\newlabel{lemma-intersection-pairing@cref}{{[originallemma][0][]0C64}{[1][0][]source}{}{}{}}
\makeatother
\crefname{originallemma}{original lemma}{original lemmas}
\Crefname{originallemma}{Original lemma}{Original lemmas}
\makeatletter
\newlabel{lemma-regular-model-connected}{{0C62}{source}{}{sourcepointer-4}{}}
\newlabel{lemma-regular-model-connected@cref}{{[originallemma][0][]0C62}{[1][0][]source}{}{}{}}
\makeatother
\crefname{originallemma}{original lemma}{original lemmas}
\Crefname{originallemma}{Original lemma}{Original lemmas}
\makeatletter
\newlabel{lemma-recurring-symmetric-real}{{0C5X}{source}{}{sourcepointer-5}{}}
\newlabel{lemma-recurring-symmetric-real@cref}{{[originallemma][0][]0C5X}{[1][0][]source}{}{}{}}
\makeatother
\title{The intersection form on the components of the special fibre is symmetric, negative semidefinite on the span of the fibre with kernel the multiples of the fibre, and vanishes on the fibre}
\author{The Stacks Project Authors}
\date{stacks-project:models.tex:local-snapshot-20261012; GNU FDL 1.2 or later}
\begin{document}
\maketitle
\noindent\textit{Editorial note.} The counted unit is the Stacks Project Lemma lemma-properties-form at source lines 4755-4769 together with its proof at source lines 4771-4790, a straight adjacent pairing in the chapter file. The chapter is the upstream file models.tex; the pool directory carries the longer name stacks-models and is not used as the chapter name anywhere in this card.\par\medskip
\noindent\textit{Editorial note.} One context span carries the regular-model situation at lines 4628-4637, which the retained statement names; the situation carries a reference of its own to the divisor lemma, which is therefore declared as a pointer as well, while the situation itself is not declared as a pointer so that no label is defined twice.\par\medskip
\noindent\textit{Editorial note.} This lemma is the one that several earlier cards of this lane point at by its tag 0C66, notably the cards on the weighted-intersection bound, on the numerical type of a model and on the genus nonnegative; the pointers and this counted unit are therefore consistent.\par\medskip
\noindent\textit{Editorial note.} The retained proof span starts at line 4771, the line of the opening marker read from the source, and contains exactly one proof environment.\par\medskip
\noindent\textit{Editorial note.} No citation occurs in the retained spans, so no bibliography entry is supplied.\par\medskip
\noindent\textit{Allowed standard prerequisites (editorial).}\begin{itemize}
\item The regular-model situation, the degree pairing, and the rank statement for recurring symmetric matrices. The situation fixes a discrete valuation ring, a smooth projective curve over its fraction field, a regular proper model and the components of the special fibre with their multiplicities; the degree pairing on combinations of components is symmetric and its value on two components is the degree of the restriction of the invertible sheaf attached to one of them to the other, which is nonnegative and vanishes exactly when the components are disjoint. The special fibre is connected and is a principal divisor on the model, so its pairing with every component vanishes, and the chapter's rank statement for symmetric recurring real matrices gives the semidefiniteness with kernel spanned by the multiplicity vector. The author's own lemmas of the chapter are the allowed inputs.. The counted lemma records the formal properties of the intersection form that every later computation of the chapter uses, in particular the semidefiniteness which forces the proportionality of a combination with vanishing self-intersection to the fibre.
\end{itemize}\medskip
\ifdefined\hypertarget\hypertarget{sourcepointer-1}{}\fi
\noindent\textit{Source reference.} Original equation 0C65: focused-seeds/AGP/sources/stacks-snapshot-20261011/tags.txt, source lines 15285--15285. Navigation pointer to the chapter's displayed definition of the intersection form, named by the retained statement; the cited labelled equation is an allowed input and its text is not reproduced here. The number is the stable Stacks tag of the target, looked up in the shared tag index.\par\medskip
\ifdefined\hypertarget\hypertarget{sourcepointer-2}{}\fi
\noindent\textit{Source reference.} Original lemma 0C5Z: focused-seeds/AGP/sources/stacks-snapshot-20261011/tags.txt, source lines 15279--15279. Navigation pointer carried by the retained situation and invoked by the retained proof for the vanishing of the pairing of the fibre with every component; that lemma is counted by a separate card of this lane and its interval is not reproduced here. The number is the stable Stacks tag of the target, looked up in the shared tag index.\par\medskip
\ifdefined\hypertarget\hypertarget{sourcepointer-3}{}\fi
\noindent\textit{Source reference.} Original lemma 0C64: focused-seeds/AGP/sources/stacks-snapshot-20261011/tags.txt, source lines 15284--15284. Navigation pointer invoked by the retained proof for the symmetry of the degree pairing, which is the first property of the statement; that lemma is counted by a separate card of this lane and its interval is not reproduced here. The number is the stable Stacks tag of the target, looked up in the shared tag index.\par\medskip
\ifdefined\hypertarget\hypertarget{sourcepointer-4}{}\fi
\noindent\textit{Source reference.} Original lemma 0C62: focused-seeds/AGP/sources/stacks-snapshot-20261011/tags.txt, source lines 15282--15282. Navigation pointer invoked by the retained proof for the connectedness of the special fibre, used in the analysis of the kernel; that lemma is counted by a separate card of this lane and its interval is not reproduced here. The number is the stable Stacks tag of the target, looked up in the shared tag index.\par\medskip
\ifdefined\hypertarget\hypertarget{sourcepointer-5}{}\fi
\noindent\textit{Source reference.} Original lemma 0C5X: focused-seeds/AGP/sources/stacks-snapshot-20261011/tags.txt, source lines 15277--15277. Navigation pointer invoked by the retained proof for the semidefiniteness of the form and the identification of its kernel with the multiples of the multiplicity vector; that lemma is counted by a separate card of this lane and its interval is not reproduced here. The number is the stable Stacks tag of the target, looked up in the shared tag index.\par\medskip
\begin{situation}
\label{situation-regular-model}
Let $R$ be a discrete valuation ring with fraction field $K$,
residue field $k$, and uniformizer $\pi$.
Let $C$ be a smooth projective curve over $K$ with $H^0(C, \mathcal{O}_C) = K$.
Let $X$ be a regular proper model of $C$.
Let $C_1, \ldots, C_n$ be the irreducible components of the special
fibre $X_k$. Write $X_k = \sum m_i C_i$ as in
Lemma \ref{lemma-divisor-special-fiber}.
\end{situation}

\begin{lemma}
\label{lemma-properties-form}
In Situation \ref{situation-regular-model} the symmetric bilinear form
(\ref{equation-form}) has the following properties
\begin{enumerate}
\item $(C_i \cdot C_j) \geq 0$ if $i \not = j$ with equality if and only
if $C_i \cap C_j = \emptyset$,
\item $(\sum m_i C_i \cdot C_j) = 0$,
\item there is no nonempty proper subset $I \subset \{1, \ldots, n\}$
such that $(C_i \cdot C_j) = 0$ for $i \in I$, $j \not \in I$.
\item $(\sum a_i C_i \cdot \sum a_i C_i) \leq 0$ with equality if and
only if there exists a $q \in \mathbf{Q}$ such that $a_i = qm_i$
for $i = 1, \ldots, n$,
\end{enumerate}
\end{lemma}

\begin{proof}
In the proof of Lemma \ref{lemma-intersection-pairing} we saw that
$(C_i \cdot C_j) = \deg(C_i \cap C_j)$ if $i \not = j$. This is
$\geq 0$ and $> 0 $ if and only if $C_i \cap C_j \not = \emptyset$.
This proves (1).

\medskip\noindent
Proof of (2). This is true because by Lemma \ref{lemma-divisor-special-fiber}
the invertible sheaf associated to $\sum m_i C_i$
is trivial and the trivial sheaf has degree zero.

\medskip\noindent
Proof of (3). This is expressing the fact that $X_k$ is connected
(Lemma \ref{lemma-regular-model-connected})
via the description of the intersection products given in the proof of (1).

\medskip\noindent
Part (4) follows from (1), (2), and (3) by
Lemma \ref{lemma-recurring-symmetric-real}.
\end{proof}

\end{document}
